\label{chap:literature_review}

\section{Overview of Related Work}
Managing blood glucose using personalized digital tools connects several areas of study, including physiology, sensor data processing, deep learning, and gut microbiome research. In this chapter, we review existing approaches, explain how they work, and discuss where they fall short in real-world diabetes management.

\section{Physiological Compartmental Models}
Historically, mathematical modeling of the glucose-insulin regulatory system has been dominated by compartmental differential equation models. The earliest foundational formulation was the \textbf{Bergman Minimal Model} (1979), which employs a two-compartment system of non-linear ordinary differential equations (ODEs) to capture insulin sensitivity ($S_I$) and glucose effectiveness ($S_G$):
\begin{align}
    \frac{dG(t)}{dt} &= -(p_1 + X(t))G(t) + p_1 G_b \\
    \frac{dX(t)}{dt} &= -p_2 X(t) + p_3 (I(t) - I_b)
\end{align}
where $G(t)$ denotes plasma glucose, $I(t)$ denotes plasma insulin, $X(t)$ represents insulin action in the interstitial space, and $p_1, p_2, p_3$ are patient-specific kinetic constants.

To capture the complex multi-organ dynamics of meal digestion and insulin kinetics, the \textbf{UVA/Padova Type 1 Diabetes Simulator} was developed by Cobelli, Dalla Man, and colleagues \cite{cobelli2011uva}. The UVA/Padova simulator expands the minimal model into a 13-compartment non-linear ODE system tracking gastric emptying, two-compartment gut glucose absorption, hepatic glucose production (HGP), renal clearance, and peripheral glucose utilization. While the UVA/Padova system achieved regulatory approval from the United States Food and Drug Administration (FDA) as an in silico substitute for preclinical animal trials in artificial pancreas algorithms, its clinical deployment for free-living Type 2 Diabetes management faces insurmountable barriers:
\begin{itemize}
    \item \textbf{Intractable Parameter Identification:} Identifying patient-specific constants requires invasive, multi-hour triple-tracer intravenous glucose tolerance tests (IVGTT) or hyperinsulinemic-euglycemic clamps, which are clinically infeasible for consumer digital twins.
    \item \textbf{Omission of Heterogeneous Modalities:} Compartmental ODEs do not possess mechanisms to incorporate real-time fitness wearable streams (variable METs and heart rates), dietary macronutrient splits (fiber, fats, protein interactions), or metagenomic gut profiles.
\end{itemize}

\section{Machine Learning Approaches for Glycemic Forecasting}
With the advent of Continuous Glucose Monitoring (CGM) sensors, researchers shifted toward data-driven statistical and machine learning algorithms.

\subsection{Autoregressive and Kernel Methods}
Early work on glucose forecasting relied on standard time-series techniques like Autoregressive Integrated Moving Average (ARIMA), Auto-Regressive with Exogenous Inputs (ARX), and Support Vector Regression (SVR) \cite{georga2013evaluating}. For instance, Georga et al. \cite{georga2013evaluating} used an SVR model with radial basis function (RBF) kernels to forecast glucose over short 15- to 30-minute horizons. These models can follow stable fasting glucose trends during the night, but they struggle after meals. Because they do not take meal macronutrients into account, they cannot anticipate the sudden rise in blood sugar following food intake, resulting in predictions that lag behind real physiological spikes by 20 to 40 minutes.

\subsection{Tabular Gradient-Boosted Decision Trees}
To include nutritional values and patient blood markers, later studies turned to tree-based models like XGBoost, LightGBM, and CatBoost. Zeevi et al. \cite{zeevi2015personalized} tracked 800 participants over roughly 46,000 meals and used gradient-boosted trees to predict post-meal incremental Area Under the Curve (iAUC). More recently, Mendes et al. \cite{mendes2024cgmacros} introduced the CGMacros dataset and tested tabular regression models using wearable and CGM data.

While tabular methods handle static features well, they have a key limitation: they reduce continuous time-series curves into single summary numbers, such as peak glucose ($G_{\max}$) or 2-hour iAUC. Because of this, they cannot output the continuous 120-minute curve ($\mathbf{Y} \in \mathbb{R}^{24}$) needed to see the full post-meal trajectory and rate of change.

\section{Deep Sequential Architectures and Recurrent Networks}
To work directly with continuous time-series data without flattening it, deep learning models have become popular in diabetes research \cite{zhu2021deep}. Recurrent Neural Networks (RNNs), especially Long Short-Term Memory (LSTM) \cite{mirshekarian2019using} and Gated Recurrent Units (GRU) \cite{martinsson2020blood}, use internal gates to pass information along the sequence:
\begin{align}
    \mathbf{f}_t &= \sigma(\mathbf{W}_f [\mathbf{h}_{t-1}, \mathbf{x}_t] + \mathbf{b}_f) \\
    \mathbf{i}_t &= \sigma(\mathbf{W}_i [\mathbf{h}_{t-1}, \mathbf{x}_t] + \mathbf{b}_i) \\
    \mathbf{c}_t &= \mathbf{f}_t \odot \mathbf{c}_{t-1} + \mathbf{i}_t \odot \tanh(\mathbf{W}_c [\mathbf{h}_{t-1}, \mathbf{x}_t] + \mathbf{b}_c) \\
    \mathbf{h}_t &= \sigma(\mathbf{W}_o [\mathbf{h}_{t-1}, \mathbf{x}_t] + \mathbf{b}_o) \odot \tanh(\mathbf{c}_t)
\end{align}

Even though recurrent networks are designed for time-series, they face two main problems when predicting post-meal glucose trajectories:
\begin{enumerate}
    \item \textbf{Spike Flattening and Averaging Out:} When forecasting multiple steps ahead (120 minutes or 24 steps), recurrent models tend to lose sharp details. Because standard training uses Mean Squared Error (MSE), the model learns that outputting a flat curve near the average value ($\approx 120\text{ mg/dL}$) yields a lower loss, which causes it to miss the real height of post-meal spikes.
    \item \textbf{Difficulty Retaining Static Clinical Context:} In typical LSTMs, static features like lab tests, age, and gut microbiome data are either copied across every time step or fed only at the start. Over long sequences, the dynamic inputs overwrite these static values, making the model lose track of individual patient traits.
\end{enumerate}

\section{Attention Mechanisms and Transformers in Healthcare Time-Series}
The Transformer architecture introduced by Vaswani et al. \cite{vaswani2017attention} replaced step-by-step recurrence with Multi-Head Self-Attention (MHSA). Instead of compressing history into a single hidden vector, self-attention lets the network compare any two points in time directly, assigning attention weights based on relevance:
\begin{equation}
    \text{Attention}(\mathbf{Q}, \mathbf{K}, \mathbf{V}) = \text{softmax}\left(\frac{\mathbf{Q} \mathbf{K}^T}{\sqrt{d_k}}\right) \mathbf{V}
\end{equation}

For multi-horizon forecasting, Lim et al. \cite{lim2021temporal} proposed the Temporal Fusion Transformer (TFT). TFT introduced Gated Residual Networks (GRN) to process static variables and filter out unhelpful features while keeping key patient signals. In this work, we build on these ideas by combining self-attention for pre-meal sensor trends with gated residual units for patient metadata, allowing the model to look at both continuous glucose momentum and clinical traits together.

\section{Gut Metagenomics in Metabolic Regulation and Representation Learning}
The human gastrointestinal tract harbors trillions of microorganisms whose collective genome (the gut metagenome) encodes metabolic pathways that directly influence host glucose homeostasis \cite{turnbaugh2006obesity}. Microbial fermentation of non-digestible polysaccharides yields Short-Chain Fatty Acids (SCFAs)—predominantly acetate, propionate, and butyrate \cite{canfora2015short}. Butyrate binds to G-protein-coupled receptors (GPR41/FFAR3 and GPR43/FFAR2) on intestinal enteroendocrine L-cells, stimulating the secretion of glucagon-like peptide-1 (GLP-1) and peptide YY (PYY), which enhance glucose-dependent insulin secretion and delay gastric emptying \cite{canfora2015short}. Conversely, pro-inflammatory lipopolysaccharide (LPS) shed by Gram-negative enterobacteria triggers metabolic endotoxemia and systemic insulin resistance \cite{asnicar2021microbiome}.

\subsection{Limitations of Traditional Dimensionality Reduction (PCA)}
Gut microbiome profiling records the presence or absence of thousands of bacterial species ($D \ge 1,000$), while study groups usually have fewer than 100 participants ($N < 100$). Because there are far more features than samples ($D \gg N$), we have to reduce the number of features so the model does not overfit.

Earlier studies, such as Zeevi et al. \cite{zeevi2015personalized}, used linear Principal Component Analysis (PCA) or Principal Coordinate Analysis (PCoA). However, linear PCA has two notable limitations for this kind of data:
\begin{enumerate}
    \item \textbf{Mismatch with Data Distribution:} PCA assumes continuous data that follows a normal distribution. Microbiome sequencing tables, however, consist of sparse zero-and-one presence-absence indicators ($x_j \in \{0, 1\}$), which violates this assumption.
    \item \textbf{Inability to Capture Non-Linear Microbial Consortia:} Bacteria function within non-linear ecological guilds (e.g., cross-feeding networks where Species A produces lactate which is metabolized by Species B into butyrate). Linear projections fail to model non-linear co-occurrence structures.
\end{enumerate}

\subsection{Variational Autoencoders ($\beta$-VAEs)}
To learn smooth, continuous, non-linear representations of high-dimensional data, Kingma and Welling \cite{kingma2013auto} introduced the \textbf{Variational Autoencoder (VAE)}. Rather than mapping inputs to deterministic point vectors, a VAE maps an input $\mathbf{x}$ to the parameters of a latent probability distribution $q_\phi(\mathbf{z}|\mathbf{x}) = \mathcal{N}(\boldsymbol{\mu}, \boldsymbol{\Sigma})$, optimizing the Evidence Lower Bound (ELBO):
\begin{equation}
    \mathcal{L}_{\text{ELBO}}(\theta, \phi; \mathbf{x}) = \mathbb{E}_{q_\phi(\mathbf{z}|\mathbf{x})}\big[\log p_\theta(\mathbf{x}|\mathbf{z})\big] - \mathcal{D}_{\text{KL}}\big(q_\phi(\mathbf{z}|\mathbf{x}) \,\|\, p(\mathbf{z})\big)
\end{equation}

Higgins et al. \cite{higgins2017beta} introduced the \textbf{$\beta$-VAE}, which introduces an adjustable hyperparameter $\beta > 0$ scaling the Kullback-Leibler (KL) divergence penalty. By modulating $\beta$, the network balances reconstruction fidelity against latent channel capacity, encouraging the discovery of disentangled, statistically independent latent factors. In this thesis, we customize the $\beta$-VAE architecture with stochastic taxa-masking denoising and $\beta$-annealing to successfully encode 1,979 binary gut species into an 8-dimensional latent manifold without overfitting.

\section{Digital Twins and Causal Simulation in Metabolic Care}
A Digital Twin is a computational model of a person's physiology that updates using real-time sensor streams and can simulate future outcomes \cite{voigt2021digital}. In metabolic health, Shamanna et al. \cite{shamanna2020twin} tested a digital twin platform in a clinical trial and found that personalized twin-based guidance helped reduce HbA1c and improved metabolic control in many patients.

Most current digital twin platforms, however, rely on static if-then rules. For a digital twin to give practical advice, it needs to support causal "what-if" simulations \cite{pearl2009causality}. By simulating questions like *"What would happen to this patient's blood sugar if they reduced carbs by 50\% or took a 30-minute walk after lunch?"*, the system allows patients to test practical options before eating.

\section{Summary of Gaps and Proposed GlycoTwin Solutions}
Table \ref{tab:literature_comparison} synthesizes the key limitations of existing paradigms and highlights the solutions introduced by the GlycoTwin framework in this thesis.

\begin{table}[!htbp]
\centering
\caption{Comprehensive Comparison: Existing Paradigms vs. GlycoTwin Framework}
\label{tab:literature_comparison}
\resizebox{\textwidth}{!}{%
\begin{tabular}{lp{3.5cm}p{3.8cm}p{5.5cm}}
\toprule
\textbf{Capability / Dimension} & \textbf{Physiological Models (UVA/Padova)} & \textbf{Tabular ML (Zeevi / CGMacros)} & \textbf{GlycoTwin Framework (Ours)} \\
\midrule
\textbf{Multi-Modal Telemetry} & None (Synthetic blood glucose only) & Tabular features only (no continuous sequence modeling) & Synchronized dual CGM (1-min), Fitbit HR/METs, Diet, Blood, Gut Microbiome. \\
\midrule
\textbf{Prediction Horizon} & Short-term continuous ($<30$ min ahead) & Scalar only ($G_{\max}$ or iAUC summary statistic) & Continuous 120-minute multi-horizon postprandial trajectory ($\mathbf{Y} \in \mathbb{R}^{24}$). \\
\midrule
\textbf{Microbiome Encoding} & None & Linear PCA / Bray-Curtis distance & Deep $\beta$-VAE with stochastic taxa-masking denoising (1.0000 ROC-AUC). \\
\midrule
\textbf{Neural Architecture} & Mechanistic ODEs (requires invasive IVGTT) & Gradient Boosted Trees (XGBoost, LightGBM) & TransformerForecaster (MHSA + Gated Residual Networks + Residual Delta). \\
\midrule
\textbf{Spike Fidelity} & Highly variable & Mean-biased scalar & Composite Excursion Loss ($\mathcal{L}_{\text{MSE}} + \text{Peak} + \text{Velocity}$). \\
\midrule
\textbf{Clinical Safety} & Not certified on free-living cohorts & Unassessed on continuous trajectories & ISO 15197 / Clarke Error Grid Certified (\textbf{98.47\% Zone A+B}). \\
\midrule
\textbf{Intervention Engine} & Closed-loop insulin pump dosing only & Retrospective statistical correlation & Real-Time Causal "What-If" Simulator \& Mobile PWA. \\
\bottomrule
\end{tabular}%
}
\end{table}
